\subsection{Definición estructural y función de las constantes de Planck}
\label{sec:ii-definicion-planck}

\begin{definition}[Determinación estructural de la acción angular]
En la recta orientada \(\mathcal L_S\), la determinación HMT de la
acción angular es el par de secciones
\(\hbar_{\mathrm{pre}},\hbar_{\mathrm{ret}}\) de
\eqref{eq:el-secciones-accion}, junto con su identificación por retorno.
Su estructura reúne la valoración de la norma local
\(\varphi\mathrm N_{G_H}(\alpha)\), el carácter \(H_5\) y la
continuación regional \(\mathcal C_\pi\). La acción por vuelta completa
es \(h_a=2\pi\hbar_a\) en cada orientación \(a\).
\end{definition}

Esta definición conserva tres determinaciones sucesivas. El cociente
local de \eqref{eq:el-h4} fija las dieciséis hojas de la norma; la
desigualdad \eqref{eq:el-decada} fija su orden decimal. El carácter y
la renovación normalizada de \eqref{eq:revision-renovacion} determinan
después las coordenadas anterior y posterior. La unicidad de la cola
utiliza su primer peso y su regla de concatenación; la distribución de
los coeficientes de \(H_5\) conserva el carácter declarado en la
sección~\ref{sec:revision-planck}. En la rama positiva, la información
multiplicativa del retorno es el invariante
\begin{equation}
 \frac{\hbar_{\mathrm{pre}}}{\hbar_{\mathrm{ret}}}
 =R_{\mathrm{act}}=\frac{H_5}{\eta_{\mathrm{ret}}}>1.
 \label{eq:ii-definicion-planck-retorno}
\end{equation}

Su función física se expresa mediante la contracción entre acción y
frecuencia: una sección \(\hbar_a\) asigna energía a una frecuencia
angular, mientras \(h_a\) expresa la acción de una vuelta. El cociente
\eqref{eq:ii-definicion-planck-retorno} conserva la modificación
estructural entre ambas lecturas aun cuando se cambia la base común
de unidades. La identificación con la constante física y su precisión
se examinan en la carta de la sección~\ref{sec:revision-planck-contraste},
con las dos coordenadas de retorno diferenciadas.
